\subsection{FRM Exam 2008: which stress testing statements are true?}

\textbf{Enunciado}

\emph{Which of the following statements about stress testing are true?}
\begin{enumerate}[label=\Roman*.]
\item Stress testing can complement VAR estimation in helping risk managers identify crucial vulnerabilities in a portfolio.
\item Stress testing allows users to include scenarios that did not occur in the lookback horizon of the VAR data but are nonetheless possible.
\item A drawback of stress testing is that it is highly subjective.
\item The inclusion of a large number of scenarios helps management better understand the risk exposure of a portfolio.
\end{enumerate}
\begin{enumerate}[label=\alph*.]
\item I and II only.
\item III and IV only.
\item I, II, and III only.
\item \hl{I, II, III, and IV.}
\end{enumerate}

\subsubsection*{Justificación}

(I) y (II) son verdaderas (stress testing complementa VaR y permite escenarios no observados). (III) es verdadera (selección de escenarios es subjetiva). (IV) es verdadera en sentido operativo: más escenarios, si están bien estructurados, amplían la comprensión del riesgo.











